\documentclass[11pt]{article}
\usepackage[margin=1in]{geometry}
\usepackage{amsmath,amssymb,amsthm,mathtools}
\usepackage{booktabs}
\usepackage{enumitem}
\usepackage{microtype}
\usepackage[hidelinks]{hyperref}

\newcommand{\A}{\mathcal A}
\newcommand{\K}{\mathcal K}
\newcommand{\R}{\mathbb R}
\newcommand{\T}{\mathbb T}
\newcommand{\dd}{\mathrm d}
\newcommand{\In}{I_n}
\newcommand{\Rn}{R_n}
\newcommand{\Un}{U_n}
\newcommand{\Sn}{S}
\newtheorem{definition}{Definition}
\newtheorem{proposition}{Proposition}
\newtheorem{theorem}{Theorem}
\newtheorem{lemma}{Lemma}

\title{V22 P3--T03: A Controlled Source-Chart Refinement Limit\\
for the Six-Channel Transport Candidate}
\author{The \AE ther-Flow Research Project}
\date{2026-08-09}

\begin{document}
\maketitle

\begin{abstract}
This draft/control manuscript constructs one finite-to-continuum refinement
family for the proposal-only P3--T02 six-channel source transport law.  The
family uses directed source-chart complexes, exact point sampling, a declared
cellwise-constant comparison map, and a monotone upwind evolution under an
explicit source CFL condition.  We prove contraction, reconstruction,
consistency, global evolution-error, boundary, and fixed-mode
principal-factor bounds.  Uniform, smoothly graded, and source-indicator-
adapted computations have nonzero errors that decrease under refinement; an
independent implementation and an inflow-boundary study agree.  These are
source-chart numerical statements, not target-geometry interpolation.  The
P3--T02 candidate supplies no quadratic physical cone or Lorentzian signature,
so their preservation status is precisely unresolved rather than passed.  No
ontology, source law, physical geometry, or Distance-to-GR status is promoted.
\end{abstract}

\section{Inherited candidate and claim boundary}

On a source chart with labels $(s^0,s^1,s^2,s^3)$, the P3--T02 fixture uses
$\tau=s^0$ and six fixed source vectors
\begin{equation}
  X_A=\partial_0+v_A^j\partial_j,
  \qquad
  v_A\in\{(1,0,0),(0,1,0),(0,0,1),(1,1,0),(0,1,2),(2,0,1)\}.
  \label{eq:velocities}
\end{equation}
The continuum equation for channel $A$ is
\begin{equation}
  \partial_0 q^A+v_A^j\partial_j q^A=0,
  \label{eq:transport}
\end{equation}
with no sum over $A$.  The coefficient datum remains proposal-only.  The
chart is a computational source presentation, not a target atlas or physical
spacetime coordinate system.

For the bounded comparison fixture, identify the three spatial source labels
periodically and write $\T^3_{\rm src}=\R^3/\mathbb Z^3$.  This quotient is a
test domain chosen before any target comparison.  It is not adopted as
$M_{\rm src}$ and carries no physical metric or measure.  The exact continuum
limit object is the six-component translation semigroup
\begin{equation}
  (S_A(t)f)(s)=f(s-v_A t),\qquad t\ge0,
  \label{eq:semigroup}
\end{equation}
on periodic source scalars.  Equation \eqref{eq:semigroup}, not agreement with
GR, defines the limit.

\section{The refinement family and exact comparison maps}

\begin{definition}[Directed product complexes]
For each level $n$, choose three cyclic source-coordinate partitions
$0=s_{j,0}<\cdots<s_{j,N_{j,n}-1}<1$.  Let $\K_n$ be their directed product
cell complex, let
\begin{equation}
  h_{j,i}=s_{j,i}-s_{j,i-1}\pmod 1,
  \qquad h_n=\max_{j,i}h_{j,i},
  \label{eq:mesh}
\end{equation}
and require $h_n\to0$.  For the irregular and adapted families, also require
a fixed mesh-ratio bound $h_n/h_{n,\min}\le C_{\rm mesh}$.
\end{definition}

Three preregistered variants are used:
\begin{enumerate}[leftmargin=*]
  \item uniform partitions;
  \item smoothly graded partitions
  $s_{j,i}=u_i+a_j\sin(2\pi u_i)/(2\pi)$ with
  $(a_1,a_2,a_3)=(0.22,-0.17,0.13)$; and
  \item source-indicator-adapted partitions with fixed source phases and
  positive cell weights proportional to
  \begin{equation*}
    \left[1+0.75\left|\sin\bigl(2\pi(u_{i+1/2}+\phi_j)\bigr)\right|\right]^{-1}.
  \end{equation*}
\end{enumerate}
All rules are frozen before target access.  The third case adapts to a declared
source variation indicator, not to target error, and is fixed during each run.

The exact sampling and reconstruction maps are
\begin{align}
  I_n &: C(\T^3_{\rm src})\longrightarrow\R^{|\K_n^0|},
       &(I_nf)_{i_1i_2i_3}&=f(s_{1,i_1},s_{2,i_2},s_{3,i_3}),
       \label{eq:sampling}\\
  R_n &: \R^{|\K_n^0|}\longrightarrow L^\infty(\T^3_{\rm src}),
       &(R_nz)|_{C_{i_1i_2i_3}}&=z_{i_1i_2i_3}.
       \label{eq:reconstruction}
\end{align}
Here $C_{i_1i_2i_3}$ is the directed half-open source cell whose lower node is
the listed node.  No target object is interpolated.  The norms
$\|z\|_{\ell^\infty}=\max_i|z_i|$ and
$\|f\|_{\infty,\mathrm{src}}=\sup_s|f(s)|$ are scalar comparison norms;
they are not physical distance or field norms.

\begin{lemma}[Reconstruction bound]
For $f\in C^1(\T^3_{\rm src})$,
\begin{equation}
 \|R_nI_nf-f\|_{\infty,\mathrm{src}}
 \le \sum_{j=1}^3h_{j,\max}\|\partial_jf\|_{\infty,\mathrm{src}}
 \le 3h_n\max_j\|\partial_jf\|_{\infty,\mathrm{src}}.
 \label{eq:reconstructionbound}
\end{equation}
\end{lemma}

\begin{proof}
Join a point in a cell to its lower node one coordinate at a time and apply the
one-dimensional mean-value theorem.  Only source-coordinate differences are
used.
\end{proof}

\section{Finite evolution, stability, and global error}

All velocities in \eqref{eq:velocities} have nonnegative source components.
Define the one-step upwind map for channel $A$ by
\begin{equation}
 (U_{A,n}z)_i=z_i-\Delta t_n\sum_{j=1}^3
 v_A^j\frac{z_i-z_{i-e_j}}{h_{j,i_j}}.
 \label{eq:upwind}
\end{equation}
The declared local source CFL condition is
\begin{equation}
 C_{A,n}(i):=\Delta t_n\sum_j\frac{v_A^j}{h_{j,i_j}}\le1.
 \label{eq:cfl}
\end{equation}
The computations use the stricter maximum $C_{A,n}\le0.45$ and choose an
integer step count so the final comparison time is exact.

\begin{proposition}[Monotonicity and stability]
Under \eqref{eq:cfl}, $U_{A,n}$ is a convex-combination map.  It preserves
constants and scalar order, maps the input range into itself, and is an
$\ell^\infty$ contraction:
\begin{equation}
  \|U_{A,n}z-U_{A,n}w\|_{\ell^\infty}
  \le\|z-w\|_{\ell^\infty}.
  \label{eq:stability}
\end{equation}
\end{proposition}

\begin{proof}
The center weight is $1-C_{A,n}(i)\ge0$, each backward-neighbor weight is
nonnegative, and all weights sum to one.  The claims follow pointwise.
\end{proof}

\begin{theorem}[Controlled source-semigroup convergence]
Let $f\in C^2(\T^3_{\rm src})$, let $t_m=m\Delta t_n\le T$, and suppose
\eqref{eq:cfl} holds.  Then
\begin{align}
 &\|U_{A,n}^{m}I_nf-I_nS_A(t_m)f\|_{\ell^\infty}\nonumber\\
 &\quad\le\frac{T}{2}\left[
 \Delta t_n\|(v_A\cdot\nabla)^2f\|_{\infty,\mathrm{src}}
 +\sum_{j=1}^3v_A^jh_{j,\max}
 \|\partial_j^2f\|_{\infty,\mathrm{src}}\right].
 \label{eq:globalerror}
\end{align}
Consequently, if $h_n\to0$ and $\Delta t_n=O(h_n)$, then
$R_nU_{A,n}^mI_nf\to S_A(t_m)f$ in the declared source sup norm, with the
additional reconstruction error \eqref{eq:reconstructionbound}.
\end{theorem}

\begin{proof}
Taylor expansion in source time gives a one-step temporal remainder bounded by
$\Delta t_n^2\|(v_A\cdot\nabla)^2f\|/2$.  The backward difference in direction
$j$ differs from $\partial_jf$ by at most
$h_{j,\max}\|\partial_j^2f\|/2$.  Insert the exact solution into
\eqref{eq:upwind}, sum the local residual over at most $T/\Delta t_n$ steps,
and use the contraction \eqref{eq:stability}.  Finally apply
\eqref{eq:reconstructionbound}.
\end{proof}

The theorem controls nonzero truncation error; it does not infer a continuum
statement from an exact finite residual.  It applies to all three mesh
families because each has $h_n\to0$, a bounded mesh ratio, and the same fixed
continuum semigroup.

\section{Boundary and finite-size protocol}

For a source slab, prescribe channel $A$ only on faces with negative inward
directional derivative, equivalently the declared inflow faces for $v_A$.
No independent outflow value is accepted.  If the discrete inflow trace has
maximum error $\varepsilon_{\partial,n}$, the convex update propagates at most
$\varepsilon_{\partial,n}$ in addition to the interior consistency bound.

\begin{proposition}[Protected finite-size window]
At a point $(t,s)$ whose exact backward $v_A$ characteristic remains in the
slab or first meets a declared inflow face with supplied trace, the finite
slab comparison uses only the initial data and that trace.  Changing data on
an outflow face cannot change the accepted update because outflow data are not
arguments of the scheme.  Outside this backward-characteristic condition, no
finite-size guarantee is claimed.
\end{proposition}

The executable boundary control solves $u_t+u_x=0$ on $[0,1]$ with an exact
source inflow trace.  Its maximum errors at $N=16,24,32,48$ are respectively
$1.492\times10^{-2}$, $1.004\times10^{-2}$,
$7.602\times10^{-3}$, and $4.791\times10^{-3}$.  The nonzero decrease is the
accepted boundary evidence.

\section{Universality and independent computation}

The six-channel maximum nodal errors at the preregistered levels are:
\begin{center}
\small
\begin{tabular}{lrrrrr}
\toprule
Mesh family & $N=6$ & $N=8$ & $N=12$ & $N=16$ & terminal order\\
\midrule
uniform & 0.09034 & 0.07784 & 0.06168 & 0.04897 & 0.802\\
graded & 0.09637 & 0.08640 & 0.06964 & 0.05826 & 0.628\\
source-indicator-adapted & 0.10368 & 0.08555 & 0.07137 & 0.06333 & 0.552\\
\bottomrule
\end{tabular}
\end{center}
Every series decreases strictly and remains nonzero.  The largest observed
mesh ratio is below $1.72$, within the preregistered bound $2$.  The different
microscopic partitions therefore agree with the same analytic limit under the
declared hypotheses; this is limited universality for the numerical source
presentation, not universality of physical microphysics.

A separate one-dimensional implementation imports no primary-model code.  Its
levels and errors are
\begin{equation*}
\begin{split}
  N&=(16,24,32,48,64),\\
  e_N&=(0.01423,0.00958,0.00727,0.00488,0.00366).
\end{split}
\end{equation*}
The observed orders are $0.975$, $0.962$, $0.984$, and $0.993$.  This
independently checks the selected monotone refinement mechanism.

\section{Principal-factor comparison and property status}

On a uniform source grid, the frozen spatial factor for a fixed integer mode
$m$ is
\begin{equation}
 \Lambda_{A,n}(m)=\sum_jv_A^j\frac{1-e^{-2\pi i m_jh_n}}{h_n}.
 \label{eq:discretesymbol}
\end{equation}
It obeys
\begin{equation}
 |\Lambda_{A,n}(m)-2\pi i\,v_A\cdot m|
 \le\frac12\sum_jv_A^jh_n|2\pi m_j|^2.
 \label{eq:symbolbound}
\end{equation}
For fixed $m=(1,1,1)$, the largest factor errors at
$N=8,16,32,64$ are $7.276$, $3.685$, $1.849$, and $0.925$.  The product of
the six factors also converges by a finite telescoping-product bound.  Factor
count and the fixed P3--T02 coefficient identity are preserved.

The monotone scheme preserves constants and the range of each scalar.  In
particular, the manufactured $q^6\in[2.92,3.08]$ remains nonzero.  Preservation
of the continuum rank-two differential condition would require a separate
$C^1$ error estimate on the relevant ratios; the present sup-norm theorem does
not supply it, so that property remains unresolved for P3--T04.

Most importantly, P3--T02 constructed neither a quadratic physical causal
cone nor a Lorentzian signature.  The explicit precondition tests therefore
return
\begin{equation}
 \texttt{upstream\_object\_undefined\_not\_preserved\_or\_lost}
 \label{eq:undefinedproperties}
\end{equation}
for both objects.  Source evolution orientation and six principal factors are
not substitutes for physical causality or signature.  Reporting those
physical properties as preserved would be a hard failure.

\section{Refuter cases and scope limits}

The following variants are rejected by the result:
\begin{enumerate}[leftmargin=*]
  \item a source CFL number above one under the declared explicit scheme;
  \item a family for which $h_n$ does not tend to zero;
  \item an adaptive mesh with unbounded ratio and no replacement proof;
  \item independent values prescribed on an outflow face;
  \item a mesh, tolerance, or comparison map chosen after target-GR access;
  \item any target metric or atlas used to define interpolation or error;
  \item a zero finite residual promoted as a continuum theorem;
  \item an undeclared boundary influence entering the comparison window;
  \item a physical cone, signature, scale, or metric inferred from the six
        convergent source transport factors; and
  \item process validation treated as mathematical or physics proof.
\end{enumerate}

The source-chart family is not passive-chart invariant as a discrete object;
only the continuum P3--T02 law has the stated passive naturality.  Different
regular source charts are allowed only as separately declared comparison
families whose continuum errors are bounded.  This limitation blocks adoption
of $\K_n$ as fundamental ontology.

\section{Distance-to-GR and decisive result}

\begin{center}
\small
\begin{tabular}{p{0.29\linewidth}p{0.64\linewidth}}
\toprule
Burden & Status after this packet\\
\midrule
Source ontology primitives & Unchanged; the source arena, $\tau$, $X_A$, and $\K_n$ remain proposal-only debt.\\
Source equivalence $\mathrm{EqSrc}$ & Continuum passive naturality is inherited; discrete chart-family invariance is not established.\\
$\mathrm{RetainH}$ and $\mathrm{GenH}$ & Unchanged and not derived.\\
$\mathrm{ObsLoc}_{\mathrm{lc}}$ and $\mathrm{Resp}_{\mathrm{lc}}$ & Scalar range and $q^6$ margin are controlled; derivative-rank preservation remains unresolved.\\
$M_{\mathrm{src}}$ & No manifold adoption; the periodic source chart is a computational fixture only.\\
$g_{\mathrm{eff}}$ & Blocked; principal-factor convergence selects no quadratic cone, signature, or scale.\\
Matter coupling and Einstein equations & Blocked.\\
Finite-variation robustness & A controlled first-order source-semigroup refinement now exists for the fixed P3--T02 fixture; physical and chart-independent robustness remain open.\\
Benchmark/Gate status & Blocked; no protected decision requested.\\
Current route freeze & Not frozen: P3--T03 adds a distinct controlled-refinement payload and routes P3--T04.\\
Hard-fail status & CFL, nonrefining, unbounded-mesh, boundary, target-tuning, target-interpolation, zero-remainder, and physical-property overreads fail.\\
\bottomrule
\end{tabular}
\end{center}

The Candidate Constructor result is \emph{constructed candidate}:
`CAND-V22-B1-SIX-TRANSPORT-REFINEMENT-V1` consists of
\eqref{eq:mesh}--\eqref{eq:upwind}, the exact maps
\eqref{eq:sampling}--\eqref{eq:reconstruction}, and the limit object
\eqref{eq:semigroup}.  Under the declared smoothness, mesh, and CFL
hypotheses, the analytic bounds and independent nonzero-error computations
establish controlled convergence of this source-chart finite family.  They do
not establish a physical cone, Lorentzian signature, scale, effective metric,
matter coupling, Einstein equations, or benchmark recovery.

The next admissible V22 packet after checkpoint is P3--T04, which must use the
bounded refinement result to classify backgrounds, linearize the relevant
response, test derivative-rank adequacy, and preserve the unresolved physical
geometry boundary.

\end{document}
